\originalchapter{变分原理与非简并微扰}
\originalsection{变分上界}
许多势无法解析求解. 近似方法的目标是用已知可解模型得到有条件、可检验的估计. 若自伴$H$下界为离散基态$E_0$, 用完备本征基展开归一试探态$\psi=\sum_nc_nu_n$, 得 $\langle H\rangle=\sum_n|c_n|^2E_n\ge E_0\sum_n|c_n|^2=E_0$. 连续谱贡献也高于谱下界, 用谱测度可得到相同不等式. 因而试探态给\emph{基态能量的上界}; 试探族越小, 最优化结果未必越准确.

压窄波包会增大动量宽度和动能, 放宽会使束缚势能升高; 优化是在这两种代价之间找平衡, 不是任意把函数调得很尖. 取谐振子Hamilton算符和宽度为$a$的Gaussian试探态. 前文已算 $\langle x^2\rangle=a^2$, $\langle p^2\rangle=\hbar^2/(4a^2)$, 所以
$E(a)=\hbar^2/(8ma^2)+m\omega^2a^2/2$. 求导令零给$a^2=\hbar/(2m\omega)$, 最小值$\hbar\omega/2$. 此次恰等于精确基态, 因为试探族包含真正基态; 对一般势并不保证取等.

变分法的参数控制形状, 不是新的物理可观测量. 在实际问题中应先满足边界与归一化, 再优化; 不满足算符定义域的尖锐函数可能导致动能发散, 不能只算势能后报告一个很低的数.

\originalsection{非简并微扰展开}
定常微扰跟踪原有能级在小改变后的能量与态. $\lambda$为无量纲强度或记账参数, $W$为能量算符, 上标$(0),(1),(2)$是各阶系数, 真实能量改变量要乘$\lambda,\lambda^2$, 不是把系数当成已乘强度的位移. 设$H(\lambda)=H_0+\lambda W$, $W$厄米. 假定$E_n^{(0)}$是孤立非简并本征值, 且在所讨论$\lambda$邻域有足够光滑的本征对. 这是应用前提, 并非任意小系数都自动保证级数收敛. 写
\[
E_n=E_n^{(0)}+\lambda E_n^{(1)}+\lambda^2E_n^{(2)}+\cdots,
\quad
\ket n=\ket{n^{(0)}}+\lambda\ket{n^{(1)}}+\cdots.
\]
固定相位和尺度为中间归一化$\braket{n^{(0)}}n=1$, 因而$\braket{n^{(0)}}{n^{(1)}}=0$. 最后若计算概率须再归一化全态.

比较本征方程的一阶系数:
$(H_0-E_n^{(0)})\ket{n^{(1)}}=-(W-E_n^{(1)})\ket{n^{(0)}}$. 左乘$\bra{n^{(0)}}$给$E_n^{(1)}=W_{nn}$. 左乘不同$m$的本征bra, 得
\[
\ket{n^{(1)}}=\sum_{m\ne n}
\frac{W_{mn}}{E_n^{(0)}-E_m^{(0)}}\ket{m^{(0)}}.
\]
这里$W_{mn}=\bra{m^{(0)}}W\ket{n^{(0)}}$. 二阶本征方程完整为$(H_0-E_n^{(0)})\ket{n^{(2)}}=-W\ket{n^{(1)}}+E_n^{(1)}\ket{n^{(1)}}+E_n^{(2)}\ket{n^{(0)}}$. 左乘$\bra{n^{(0)}}$得$E_n^{(2)}=\bra{n^{(0)}}W\ket{n^{(1)}}$, 因中间归一化把另一项消去. 二阶方程左乘$\bra{n^{(0)}}$后, 中间归一化消去一阶能量乘态修正项, 得
\[
E_n^{(2)}=\sum_{m\ne n}\frac{|W_{mn}|^2}{E_n^{(0)}-E_m^{(0)}}.
\]
有限维中求和自动有限; 无限维中需要求和收敛及相应算符条件, 有连续谱时还需积分. 对基态, 所有分母负, 二阶能量修正非正. 这为计算提供符号检查.

\originalsection{扰动大小与一个算例}
重要的单通道混合诊断是$|\lambda W_{mn}/(E_n^{(0)}-E_m^{(0)})|$, 不是仅看$|\lambda|$. 多通道时还应检查$\lambda^2\sum_{m\ne n}|W_{mn}|^2/(E_n^{(0)}-E_m^{(0)})^2\ll1$或相应连续谱积分; 每个小比值不自动保证总态修正小. 当两个能级很近, 很小的矩阵元也会造成大混合, 应使用下一章的简并或近简并处理.

\begin{example}
谐振子加上弱线性势$\lambda f x$, 求基态能量到二阶并与精确解比较.
\end{example}
\begin{solution}
$f$量纲为力, $+fx$的力为$-f$, 基态的$\langle x\rangle=0$, 一阶修正为零. $x=\sqrt{\hbar/(2m\omega)}(a+a^\dagger)$只把$\ket0$连到$\ket1$. 因而二阶系数$E_0^{(2)}$为
$|f\langle1|x|0\rangle|^2/(E_0^{(0)}-E_1^{(0)})=-f^2/(2m\omega^2)$.
真实二阶能量改变量为$\lambda^2E_0^{(2)}$. 配方可得
\[
\frac12m\omega^2x^2+\lambda f x
=\frac12m\omega^2\left(x+\frac{\lambda f}{m\omega^2}\right)^2
-\frac{\lambda^2f^2}{2m\omega^2}.
\]
所以精确基态能量为$\hbar\omega/2-\lambda^2f^2/(2m\omega^2)$, 与微扰一致. 这里二阶以后能量修正为零, 因为势的变化只平移中心并加常数, 不是一般扰动都如此.
\end{solution}

\originalsection{从近似系数理解能级排斥}
扰动对角元$W_{nn}$是原态对新能量项的直接平均, 非对角元则把原来不同态连接起来. 为观察反馈, 取$H=\left(\begin{smallmatrix}E_a&\lambda w\\\lambda w^*&E_b\end{smallmatrix}\right)$,$E_b>E_a$. 精确能量$(E_a+E_b)/2\pm\sqrt{(E_b-E_a)^2/4+\lambda^2|w|^2}$, 小$|\lambda w|/(E_b-E_a)$下低能下降$-\lambda^2|w|^2/(E_b-E_a)$、高能升高同量. 态的混合幅约$\lambda w/(E_b-E_a)$, 能隙很小则这个展开失效. 基态二阶下降与变分上界并不矛盾: 微扰比较的是原$H_0$能量, 变分比较的是同一个新$H$的精确基态与试探态.

对真正非谐$V=m\omega^2x^2/2+g x^4$,$g>0$, 用第2章Gaussian宽度$a$, $\langle x^4\rangle=3a^4$, 上界$E(a)=\hbar^2/(8ma^2)+m\omega^2a^2/2+3ga^4$. 设$y=a^2$, 最佳参数由$-\hbar^2/(8my^2)+m\omega^2/2+6gy=0$确定, 左边严格递增所以唯一正根. $g$增加会缩窄最佳态, 但Gaussian不含精确四次势基态, 上界不等于精确答案. 上界保证能量方向, 不自动保证波函数误差或激发能误差同样小.
